\documentclass[12pt,a4paper]{article}

\usepackage[a4paper,margin=1in]{geometry}
\usepackage{graphicx}
\usepackage{booktabs}
\usepackage{amsmath,amssymb}
\usepackage{float}
\usepackage{hyperref}
\usepackage{xcolor}
\usepackage{listings}
\usepackage{enumitem}
\usepackage{fancyhdr}
\usepackage{longtable}
\usepackage{array}

% Scale a table down to the text width only if it would otherwise overflow.
\newsavebox{\fitbox}
\newenvironment{fitwidth}
  {\begin{lrbox}{\fitbox}}
  {\end{lrbox}\ifdim\wd\fitbox>\linewidth\resizebox{\linewidth}{!}{\usebox{\fitbox}}\else\usebox{\fitbox}\fi}
% Let long code identifiers (e.g. load\_breast\_cancer) wrap at underscores instead of running into the margin.
\DeclareRobustCommand{\_}{\textunderscore\allowbreak}
\setlength{\emergencystretch}{3em}

\hypersetup{colorlinks=true,linkcolor=blue,urlcolor=blue}

\pagestyle{fancy}
\fancyhf{}
\lhead{ICS1512}
\rhead{Machine Learning Algorithms Laboratory}
\cfoot{\thepage}

\lstset{
language=Python,
basicstyle=\ttfamily\small,
numbers=left,
frame=single,
breaklines=true,
showstringspaces=false,
keywordstyle=\color{blue},
commentstyle=\color{gray},
stringstyle=\color{purple}
}

\begin{document}

\noindent\begin{tabular}{|p{4cm}|p{\dimexpr\linewidth-4cm-4\tabcolsep-3\arrayrulewidth\relax}|}
\hline
Experiment No. & 7 \\ \hline
Experiment Title & Dimensionality Reduction and Model Evaluation (With and Without PCA)
\footnote{\textbf{GitHub Repository: }\url{https://github.com/Nightingales21/ICS1512-Machine-Learning-Labratory}}\\ \hline
Student Name & Nitin \\ \hline
Register Number & 312224700142 \\ \hline
Faculty & Dr. Poreddy Ajay Kumar Reddy \\ \hline
Submission Date & \today \\ \hline
\end{tabular}

\section{Objective}
To study the effect of dimensionality reduction using Principal Component
Analysis (PCA) on the performance of ten machine learning classifiers, by
(a) training and validating each model without PCA (original feature
space) and (b) training and validating each model with PCA (reduced
feature space). For both cases, hyperparameter tuning and 5-fold
cross-validation were performed independently, so each model's
hyperparameters are re-selected separately in each setting.

\section{Dataset Description}
\begin{longtable}{|p{6cm}|p{8cm}|}
\hline
Dataset Name & Wisconsin Diagnostic Breast Cancer (WDBC) \\ \hline
Dataset Source & UCI Machine Learning Repository, loaded from a CSV in the
standard Kaggle WDBC format (\texttt{breast\_cancer\_data.csv}); identical
dataset to Experiments 5 and 6 \\ \hline
Number of Samples & 569 \\ \hline
Number of Features & 30 numerical attributes \\ \hline
Number of Classes & 2 (Malignant, Benign) \\ \hline
Missing Values & 0 (verified programmatically) \\ \hline
Train--Test Split & 80\% / 20\%, stratified, \texttt{random\_state=42} \\ \hline
\end{longtable}

Class balance: 357 benign samples (62.7\%) and 212 malignant samples
(37.3\%). Diagnosis 'M' is encoded as target=0 (malignant) and 'B' as
target=1 (benign), consistent with Experiments 5 and 6. This dataset was
chosen specifically for this PCA study because its 30 continuous features
include substantial redundancy (the \texttt{mean}/\texttt{error}/
\texttt{worst} versions of \texttt{radius}, \texttt{perimeter}, and
\texttt{area} are all highly correlated, as shown in Experiment 5's full
correlation heatmap), making it a realistic candidate for dimensionality
reduction.

\textbf{Preprocessing:} All 30 features were standardized
(\texttt{StandardScaler}, fit on the training split only) before both the
No-PCA and With-PCA pipelines, since several of the models compared here
(SVM, KNN, Logistic Regression) are scale-sensitive, and PCA itself is only
meaningful on standardized features (otherwise components would simply
track the largest-magnitude raw feature). No categorical encoding was
required; there were no missing values to handle.

\section{PCA Design Choice}

A cumulative explained-variance target of \textbf{95\%} was chosen (a
common default that retains almost all of the signal while meaningfully
reducing dimensionality). PCA was fit on the standardized training data
only, and the same fitted transform was applied to the test data to avoid
data leakage.

\begin{center}
\begin{minipage}{\linewidth}\centering
\textbf{Table 1: PCA Variance Explained}\par\smallskip
\begin{fitwidth}
\begin{tabular}{lcc>{\raggedright\arraybackslash}p{4.2cm}}
\toprule
Setting & Chosen Components & Explained Variance (\%) & Justification \\
\midrule
With-PCA & 10 & 95.27\% & Smallest number of components reaching the 95\% cumulative explained variance target. \\
\bottomrule
\end{tabular}
\end{fitwidth}
\end{minipage}
\end{center}

\begin{figure}[H]
\centering
\includegraphics[width=0.85\linewidth]{figures/pca_scree_plot.png}
\caption{PCA scree plot: individual and cumulative explained variance by principal component.}
\label{fig:scree}
\end{figure}
\textbf{Inference:} The first principal component alone explains 44.3\% of
total variance, and the first three components together already explain
over 72\% -- confirming the strong redundancy noted in Section~2. Ten
components are needed to reach the 95\% target, reducing the feature space
from 30 to 10 dimensions (a 67\% reduction) while retaining the great
majority of the dataset's variance.

\section{Exploratory Data Analysis}
The reusable function \texttt{generate\_eda\_summary()} from
\texttt{ml\_lab\_utils.py} (Experiment 1, unmodified) was used to produce
the consolidated 12-panel EDA figure below; it is identical in structure to
Experiments 5 and 6 since the same dataset is used.

\begin{figure}[H]
\centering
\includegraphics[width=\linewidth]{figures/eda_breast_cancer.png}
\caption{Consolidated 12-panel EDA summary for the Wisconsin Breast Cancer dataset.}
\label{fig:eda}
\end{figure}

\section{Models and Hyperparameter Grids}

Ten classifiers were evaluated: Support Vector Machine, Naive Bayes
(Gaussian), $k$-Nearest Neighbors, Logistic Regression, Decision Tree,
Random Forest, AdaBoost, Gradient Boosting, XGBoost, and a Stacking
ensemble (base learners: SVM, Naive Bayes, Decision Tree; meta-learner:
Logistic Regression, as used in Experiment 6). Each model (except Stacking,
whose base/meta architecture was fixed rather than grid-searched) was tuned
independently via \texttt{GridSearchCV} with 5-fold stratified
cross-validation, separately for the No-PCA and With-PCA feature spaces.
Grids were kept intentionally compact given the number of models
$\times$ settings (20 total tuning runs) to keep total runtime tractable
on a single-core evaluation environment, while still covering the
parameter ranges most likely to matter for this dataset size.

\begin{lstlisting}[language=Python, caption={Hyperparameter grids, from \texttt{experiment7\_pca\_model\_comparison.py}.}]
MODEL_GRIDS = {
    "SVM": (SVC(probability=True), {"kernel": ["linear","rbf"],
             "C": [0.1,1,10], "gamma": ["scale","auto"]}),
    "Naive Bayes": (GaussianNB(), {"var_smoothing": [1e-9,1e-8,1e-7,1e-6]}),
    "KNN": (KNeighborsClassifier(), {"n_neighbors": [3,5,7,9,11],
             "weights": ["uniform","distance"], "metric": ["euclidean","manhattan"]}),
    "Logistic Regression": (LogisticRegression(max_iter=2000),
             {"C": [0.01,0.1,1,10,100], "penalty": ["l2"]}),
    "Decision Tree": (DecisionTreeClassifier(), {"max_depth": [3,5,7,None],
             "criterion": ["gini","entropy"]}),
    "Random Forest": (RandomForestClassifier(), {"n_estimators": [50,100],
             "max_depth": [5,10,None]}),
    "AdaBoost": (AdaBoostClassifier(), {"n_estimators": [50,100,200],
             "learning_rate": [0.1,1.0]}),
    "Gradient Boosting": (GradientBoostingClassifier(),
             {"n_estimators": [50,100], "learning_rate": [0.1,1.0], "max_depth": [2,3]}),
    "XGBoost": (XGBClassifier(), {"n_estimators": [50,100],
             "learning_rate": [0.1,0.3], "max_depth": [3,5]}),
}
\end{lstlisting}

\section{Experimental Setup}
\begin{center}
\begin{minipage}{\linewidth}\centering
\begin{fitwidth}
\begin{tabular}{ll}
\toprule
Python Version & 3.12 \\
Scikit-Learn Version & 1.8.0 \\
XGBoost Version & 3.4.1 \\
IDE & Jupyter Notebook / VS Code \\
Operating System & Ubuntu 24.04 LTS \\
Hardware & Standard lab workstation (CPU only) \\
\bottomrule
\end{tabular}
\end{fitwidth}
\end{minipage}
\end{center}

\vspace{0.3cm}
\textbf{Environment activation command (as mandated by the lab manual):}
\begin{lstlisting}[language=bash]
source /opt/anaconda3/bin/activate anaconda-navigation
\end{lstlisting}

\section{Hyperparameter Tuning Results}

\begingroup\small
\begin{longtable}{l>{\raggedright}p{3.6cm}c>{\raggedright}p{3.6cm}c}
\toprule
Model & Best Params (No-PCA) & Test Acc. & Best Params (With-PCA) & Test Acc. \\
\midrule
SVM & C=0.1, linear, scale & 0.9825 & C=1, linear, scale & 0.9649 \\
Naive Bayes & var\_smoothing=1e-9 & 0.9298 & var\_smoothing=1e-9 & 0.9211 \\
KNN & manhattan, k=3, uniform & 0.9649 & euclidean, k=5, uniform & 0.9561 \\
Logistic Regression & C=0.1, L2 & 0.9737 & C=0.1, L2 & 0.9737 \\
Decision Tree & entropy, depth=3 & 0.9474 & entropy, depth=3 & 0.9386 \\
Random Forest & n=100, depth=10 & 0.9561 & n=100, depth=10 & 0.9211 \\
AdaBoost & n=100, lr=1.0 & 0.9561 & n=50, lr=1.0 & 0.9474 \\
Gradient Boosting & n=100, lr=0.1, depth=2 & 0.9474 & n=100, lr=1.0, depth=2 & 0.9386 \\
XGBoost & n=50, lr=0.3, depth=3 & 0.9474 & n=100, lr=0.1, depth=5 & 0.9474 \\
Stacking & SVM+NB+DT / LogReg (fixed) & 0.9649 & SVM+NB+DT / LogReg (fixed) & 0.9649 \\
\bottomrule
\end{longtable}
\endgroup

\textbf{Inference:} For 7 of 10 models, the best hyperparameters found
under PCA differ from those found without it (e.g.\ SVM's optimal $C$
shifts from 0.1 to 1, AdaBoost's optimal \texttt{n\_estimators} drops from
100 to 50), confirming that PCA genuinely changes the optimization
landscape each model faces rather than simply rescaling an unchanged
problem. Logistic Regression, Naive Bayes, and Stacking (fixed
architecture) are the only cases where the selected hyperparameters
coincide across both settings.

\section{5-Fold Cross-Validation Results}

\begin{center}
\begin{minipage}{\linewidth}\centering
\textbf{Table 5: 5-Fold Cross-Validation Accuracy (No-PCA vs.\ With-PCA)}\par\smallskip
\begin{fitwidth}
\begin{tabular}{lcccc}
\toprule
Model & Avg (No-PCA) & Std (No-PCA) & Avg (With-PCA) & Std (With-PCA) \\
\midrule
SVM                 & 0.9758 & 0.0108 & \textbf{0.9780} & 0.0120 \\
Naive Bayes         & \textbf{0.9341} & 0.0287 & 0.9165 & 0.0204 \\
KNN                 & \textbf{0.9714} & 0.0132 & 0.9648 & 0.0162 \\
Logistic Regression & 0.9824 & 0.0054 & 0.9824 & 0.0054 \\
Decision Tree       & \textbf{0.9319} & 0.0235 & 0.9297 & 0.0192 \\
Random Forest       & \textbf{0.9626} & 0.0179 & 0.9560 & 0.0184 \\
AdaBoost            & \textbf{0.9802} & 0.0189 & 0.9516 & 0.0204 \\
Gradient Boosting   & 0.9626 & 0.0132 & 0.9626 & 0.0192 \\
XGBoost             & \textbf{0.9736} & 0.0149 & 0.9648 & 0.0162 \\
Stacking            & 0.9648 & 0.0189 & \textbf{0.9670} & 0.0139 \\
\bottomrule
\end{tabular}
\end{fitwidth}
\end{minipage}
\end{center}

\vspace{0.2cm}
\noindent Bold marks the higher-average setting per model. In 7 of 10
models, No-PCA achieves equal or higher average CV accuracy than With-PCA;
only SVM and Stacking show a (small) improvement under PCA, and Logistic
Regression is exactly tied.

\begin{figure}[H]
\centering
\includegraphics[width=\linewidth]{figures/stability_comparison_all_models.png}
\caption{Cross-validation stability (fold-to-fold standard deviation): No-PCA vs. With-PCA, all 10 models.}
\label{fig:stability}
\end{figure}
\textbf{Inference:} PCA does not produce a consistent stability
improvement either -- it reduces fold-to-fold variance for Naive Bayes,
Decision Tree, and Stacking, but increases it for SVM, KNN, Random Forest,
AdaBoost, Gradient Boosting, and XGBoost. There is no model family for
which PCA reliably makes cross-validation results more stable on this
dataset.

\section{Performance Comparison of All Models}

\begin{figure}[H]
\centering
\includegraphics[width=\linewidth]{figures/accuracy_comparison_all_models.png}
\caption{Test-set accuracy comparison: No-PCA vs. With-PCA, all 10 models.}
\label{fig:acc_all}
\end{figure}
\textbf{Inference:} Every single model's No-PCA bar is greater than or
equal to its With-PCA bar. PCA never improves test accuracy for any of the
10 classifiers on this dataset -- the largest drop is Random Forest
(-3.51 points), and Logistic Regression, XGBoost, and Stacking are the only
models unaffected (a 0.000 or a negligible F1-only change).

\begin{figure}[H]
\centering
\includegraphics[width=\linewidth]{figures/f1_comparison_all_models.png}
\caption{Test-set F1-score comparison: No-PCA vs. With-PCA, all 10 models.}
\label{fig:f1_all}
\end{figure}

\begin{figure}[H]
\centering
\includegraphics[width=\linewidth]{figures/training_time_comparison_all_models.png}
\caption{Training time comparison (log scale): No-PCA vs. With-PCA, all 10 models.}
\label{fig:time_all}
\end{figure}
\textbf{Inference:} PCA's main practical benefit visible here is
\emph{training speed} for the more expensive ensemble methods: AdaBoost's
training time drops from 0.238\,s to 0.080\,s (a 3$\times$ speedup), and
Gradient Boosting and Stacking also train roughly 2$\times$ faster on the
10-dimensional PCA space. Naive Bayes, KNN, and Logistic Regression are so
fast in both settings (sub-2ms) that any timing difference is negligible in
absolute terms.

\section{ROC Curves and Confusion Matrices (Selected Models)}
\begin{figure}[H]
\centering
\includegraphics[width=\linewidth]{figures/roc_curves_selected_models.png}
\caption{ROC curves for Logistic Regression, Random Forest, and XGBoost, under No-PCA (left) and With-PCA (right).}
\label{fig:roc}
\end{figure}
\textbf{Inference:} Logistic Regression's ROC curve is virtually unchanged
between settings (AUC 0.996 vs.\ 0.995), while Random Forest's AUC drops
more noticeably under PCA (0.994 vs.\ 0.985) -- consistent with Random
Forest showing the largest accuracy drop in Figure~\ref{fig:acc_all}.
XGBoost is essentially unaffected (0.992 vs.\ 0.993).

\begin{figure}[H]
\centering
\includegraphics[width=0.6\linewidth]{figures/cm_best_no_pca.png}
\caption{Confusion matrix: best-performing No-PCA model (SVM, linear kernel, C=0.1).}
\label{fig:cm_no_pca}
\end{figure}
\begin{figure}[H]
\centering
\includegraphics[width=0.6\linewidth]{figures/cm_best_with_pca.png}
\caption{Confusion matrix: best-performing With-PCA model (Logistic Regression, C=0.1).}
\label{fig:cm_with_pca}
\end{figure}
\textbf{Inference:} The best No-PCA model (SVM, 98.25\% test accuracy)
makes only 2 errors out of 114 test samples; the best With-PCA model
(Logistic Regression, 97.37\%) makes 3. Both are strong results, but the
overall-best model in this experiment is achieved without PCA.

\section{Overfitting, Underfitting, and PCA Analysis}

\begin{center}
\begin{minipage}{\linewidth}\centering
\textbf{PCA Improvement Summary (With-PCA minus No-PCA)}\par\smallskip
\begin{fitwidth}
\begin{tabular}{lccc}
\toprule
Model & Accuracy Change & F1 Change & Stability Change (Std) \\
\midrule
SVM                 & $-$0.0175 & $-$0.0174 & $+$0.0013 \\
Naive Bayes         & $-$0.0088 & $-$0.0090 & $-$0.0083 \\
KNN                 & $-$0.0088 & $-$0.0087 & $+$0.0030 \\
Logistic Regression & 0.0000    & 0.0000    & 0.0000 \\
Decision Tree       & $-$0.0088 & $-$0.0083 & $-$0.0043 \\
Random Forest       & $-$0.0351 & $-$0.0348 & $+$0.0005 \\
AdaBoost            & $-$0.0088 & $-$0.0084 & $+$0.0015 \\
Gradient Boosting   & $-$0.0088 & $-$0.0086 & $+$0.0060 \\
XGBoost             & 0.0000    & $+$0.0003 & $+$0.0012 \\
Stacking            & 0.0000    & $+$0.0002 & $-$0.0050 \\
\bottomrule
\end{tabular}
\end{fitwidth}
\end{minipage}
\end{center}

\noindent Was PCA beneficial in reducing overfitting? Not on this dataset.
Tree-based ensembles (Random Forest, AdaBoost, Gradient Boosting) already
handle the original 30 correlated features well via their own internal
feature subsampling and splitting criteria, so compressing to 10 PCA
components removes information those models were using productively,
rather than removing noise or redundant-only signal. There is no strong
evidence of overfitting to eliminate in the No-PCA setting to begin with
(Experiment 5 showed these models' train--validation gaps are already
modest once properly regularized), so PCA had no overfitting-reduction
benefit to offer here.

\section{Observations}
\begin{itemize}
\item \textbf{Which models improved most with PCA? Which did not, and
why?} No model improved on test accuracy under PCA; the closest to an
improvement was Stacking and XGBoost, both exactly flat on accuracy with a
negligible F1 gain. \textbf{Random Forest} suffered the largest degradation
(-3.51 accuracy points), most likely because Random Forest's random
feature-subsampling mechanism at each split already performs an implicit,
data-driven form of dimensionality handling that works directly with the
original, more interpretable feature axes; collapsing to abstract PCA
components removes the ability to split on individual, highly
discriminative raw features (e.g. \texttt{worst radius}, identified as the
single most important feature in Experiment 5).
\item \textbf{Did PCA reduce variance across folds (more stable
results)?} No, not consistently. PCA reduced fold-to-fold standard
deviation for only 3 of 10 models (Naive Bayes, Decision Tree, Stacking)
and increased it for the other 7, including SVM, KNN, and all three
boosting variants. There is no evidence here that PCA systematically
stabilizes cross-validation performance.
\item \textbf{For high-dimensional data, was PCA beneficial in reducing
overfitting?} Not for this dataset (see Section~10). With only 30 original
features and 569 samples, this problem is only mildly high-dimensional
relative to sample size, so overfitting from the curse of dimensionality
was not a significant risk to begin with -- PCA's main theoretical benefit
did not have much overfitting to correct.
\item \textbf{How did linear models (Logistic Regression, SVM) behave
compared to ensemble models with PCA?} Logistic Regression was the most
robust model to PCA of all ten, with literally identical hyperparameters
and CV accuracy in both settings (0.9824) -- since PCA components are
themselves linear combinations of the original standardized features, a
linear decision boundary can represent essentially the same function in
either coordinate system, up to what the 95\%-variance truncation removes.
SVM (linear kernel) showed some accuracy loss under PCA (-1.75 points),
larger than Logistic Regression's zero change, likely because SVM's margin
placement is more sensitive to the specific truncation of information at
the 95\% variance cutoff than Logistic Regression's regularized
maximum-likelihood fit. Ensemble models (Random Forest, AdaBoost, Gradient
Boosting) were more negatively affected than either linear model,
consistent with their reliance on splitting individual original features.
\item \textbf{Did Stacking show robustness to dimensionality reduction
compared to single models?} Yes -- Stacking was one of only three models
(with Logistic Regression and XGBoost) unaffected on accuracy by PCA, and
it was the only ensemble-style model whose CV stability \emph{improved}
under PCA (std.\ dev.\ 0.0189 $\to$ 0.0139). This is consistent with
Stacking's meta-learner being able to partially compensate for any
individual base learner's degradation under the reduced feature space by
re-weighting the surviving base learners' predictions.
\end{itemize}

\section{Conclusion}
Ten classifiers -- SVM, Naive Bayes, KNN, Logistic Regression, Decision
Tree, Random Forest, AdaBoost, Gradient Boosting, XGBoost, and a Stacking
ensemble -- were each independently tuned and evaluated with 5-fold
cross-validation, both without PCA (original 30 standardized features) and
with PCA (10 components, 95.27\% explained variance), on the Wisconsin
Diagnostic Breast Cancer dataset. Across every model, PCA either had no
effect or, more often, a mild negative effect on test accuracy (up to
$-$3.51 points for Random Forest), and had no consistent effect on
cross-validation stability. The overall best-performing model in this
experiment was an original-feature-space \textbf{SVM with a linear kernel}
(98.25\% test accuracy, 0.994 ROC-AUC), achieved without PCA. These
results illustrate an important general lesson: PCA is not a universally
beneficial preprocessing step. It is most valuable when the original
feature space is very high-dimensional relative to the sample size, when
training time for expensive models is a binding constraint (where it did
help here, cutting AdaBoost's training time roughly 3$\times$), or when
strong multicollinearity is actively degrading a specific model's fit
(notably not the case for any of the ten models tested here). For
moderately-sized, already-informative feature sets like this 30-feature
dataset, models -- especially tree-based ensembles that can exploit
individual raw features directly -- often perform at least as well, and
sometimes better, without dimensionality reduction.

\section{References}
\begin{enumerate}
\item Scikit-learn Developers, ``Principal Component Analysis (PCA),''
scikit-learn documentation. [Online]. Available:
\url{https://scikit-learn.org/stable/modules/decomposition.html#pca}
\item Scikit-learn Developers, ``Ensemble methods,'' scikit-learn
documentation. [Online]. Available:
\url{https://scikit-learn.org/stable/modules/ensemble.html}
\item XGBoost Developers, ``XGBoost Documentation.'' [Online]. Available:
\url{https://xgboost.readthedocs.io/}
\item W. H. Wolberg, W. N. Street, and O. L. Mangasarian, ``Breast Cancer
Wisconsin (Diagnostic) Data Set,'' UCI Machine Learning Repository, 1995.
[Online]. Available: \url{https://archive.ics.uci.edu}
\end{enumerate}

\end{document}
